\section{Paso 7: diseño de experimentos}
\label{sec:paso7-diseno-experimentos}
% Origen: plan §7 (2026-11-04). Tabla de escenarios del experimento y sensibilidad: anexo G.

\textbf{Escenarios y tipo de simulación.} Se comparan los cuatro escenarios de la tabla~\ref{tab:escenarios}:
\EO{} y \EI{} validan con una triangular (1,8; 2,4; 3,5)~s, y \ET{} y \EIT{} con una triangular
(1,5; 1,8; 2,5)~s, que toma como moda el extremo optimista del rango \parencite[Exh.~10-27]{tcqsm2013}
(anexo~\ref{anx:escenarios}). El sistema es terminal \parencite{law2015}: arranca vacío a las 07:00 y termina
cuando sale el último pasajero que ingresó antes de las 09:30, sin período de calentamiento. Cada réplica es un
día hábil real con su semilla y produce un único valor de cada medida; los intervalos se arman con esos valores,
nunca con las esperas individuales de una misma mañana, que están correlacionadas.

\textbf{Números aleatorios comunes.} Los cuatro escenarios de una réplica usan la misma semilla y el mismo día:
los mismos trenes, la misma gente con las mismas velocidades y los mismos sorteos de la validación base, con un
generador por fuente de azar. Con números comunes,
\[
\operatorname{Var}(X_{a}-X_{\EO{}}) = \operatorname{Var}(X_{a}) + \operatorname{Var}(X_{\EO{}})
  - 2\operatorname{Cov}(X_{a}, X_{\EO{}}),
\]
y la covarianza positiva (un día cargado es malo para todos los escenarios) achica el intervalo de la
diferencia.

\textbf{Comparación y control del error.} Cada alternativa $a \in \{\EI{}, \ET{}, \EIT{}\}$ se compara contra
\EO{} con el intervalo $t$ apareado. Para la réplica $j$ y la medida $X$, $Z_j = X_{a,j} - X_{\EO{},j}$, y
\[
\bar Z(n) \pm t_{n-1,\,1-\alpha_s/2}\sqrt{\frac{\sum_{j=1}^{n}\left[Z_j-\bar Z(n)\right]^2}{n(n-1)}} .
\]
Con tres primarias y tres comparaciones, Bonferroni con $\alpha = 0{,}05$ y $k = 9$ fija $\alpha_s = 0{,}05/9$,
para una confianza conjunta de al menos el 95~\%. Las secundarias (espera media, largo de cola, cola máxima, utilización, tránsito del Roca hasta la cola y
coeficiente de variación entre molinetes) son descriptivas, con intervalos individuales del 95~\%.

\textbf{Cantidad de réplicas.} Se corre un piloto de $n_0 = 10$ réplicas y, para cada primaria y comparación,
\[
n^{*}_a(\beta) = \min\left\{ i \ge n_0 : t_{i-1,\,1-\alpha_s/2}\sqrt{S^2(n_0)/i} \le \beta \right\},
\]
con $\beta = 1$~s para los percentiles y $\beta = 0{,}01$ para la proporción, fijados antes de ver los datos; se usa precisión absoluta porque la relativa se indefine si la diferencia
real es casi cero.
El diseño de referencia era de 30 réplicas por escenario (120 corridas); con unos 170~s por corrida, la
producción arranca con las $n_0 = 10$ réplicas (40 corridas, en paralelo) y se agregan días solo si $n^{*}$ lo
pide.

\textbf{Regla de decisión.} Se fija antes de correr. Si el intervalo de una primaria queda entero por debajo de cero, la alternativa
mejora esa medida con evidencia estadística, y la mejora se informa con el intervalo, no solo con la
media. Si contiene al cero, no hay evidencia de diferencia con ese $n$; si además es angosto (semiancho
$\le \beta$), cualquier diferencia es despreciable en la práctica. Si las primarias no coinciden, se
describe el compromiso y no se declara una mejora global. Entre las alternativas que mejoran, se
recomienda la de menor percentil 90 si su intervalo frente a las otras también lo sostiene. Fuera de la familia principal se plantearon corridas de
sensibilidad con los contextos optimista y pesimista (anexo~\ref{anx:sensibilidad}).
